\documentclass[11pt]{article}

\usepackage{amsmath,amssymb}
\usepackage{xurl}
\usepackage[hidelinks]{hyperref}
\setlength{\emergencystretch}{2em}

% Shared commands for notes in docs/paper-gaps/.
% Notation follows the LDT paper (references/ldt-paper/).

\newcommand{\Lean}{\textsc{Lean}}
\newcommand{\leanid}[1]{\textsf{#1}}
\newcommand{\ghissue}[1]{\href{https://github.com/LionSR/MIPStarRE/issues/#1}{GitHub issue \##1}}
\newcommand{\ghpr}[1]{\href{https://github.com/LionSR/MIPStarRE/pull/#1}{GitHub PR \##1}}

% --- Paper-compatible notation ---
\newcommand{\F}{\mathbb{F}}
\newcommand{\N}{\mathbb{N}}
\newcommand{\C}{\mathbb{C}}
\newcommand{\tr}{\operatorname{tr}}
\newcommand{\ot}{\otimes}
\newcommand{\E}{\mathop{\bf E\/}}
\newcommand{\bx}{\mathbf{x}}
\newcommand{\by}{\mathbf{y}}
\newcommand{\bu}{\mathbf{u}}
\newcommand{\bv}{\mathbf{v}}

% Approximation relation (paper uses \approx_\eps and \simeq_\zeta)
\newcommand{\approxeps}{\approx_{\varepsilon}}
\newcommand{\simgeq}{\simeq_{\zeta}}

% Polynomial spaces (paper uses \polyfunc, \polysub, \polymeas)
\newcommand{\polyfunc}[3]{\operatorname{Poly}(#1,#2,#3)}
\newcommand{\polysub}[3]{\operatorname{PolySub}(#1,#2,#3)}
\newcommand{\polymeas}[3]{\operatorname{PolyMeas}(#1,#2,#3)}


\title{The Zero-Sampling Boundary in the Main Formal Theorem}
\author{}
\date{2026-05-22}

\begin{document}
\maketitle

\section{At a glance}

\begin{itemize}
  \item \textbf{Difficulty: low as a mathematical obstruction, medium as a
  formal repair.}  The scalar obstruction is immediate from the factor \(k^2\)
  in the printed final error; deciding the replacement theorem form requires a
  source-statement judgment.
  \item \textbf{Estimated weight:} no new proof machinery is expected if the
  formal theorem is corrected to require \(k>0\).  A fully formal
  counterexample would require a small explicit projective strategy with
  nonconstant point answers in degree zero.
  \item \textbf{Mathlib/project split:} approximately \textbf{5\% Mathlib} and
  \textbf{95\% project-local}.  The only arithmetic is the simplification of the
  displayed error at \(k=0\); the issue concerns the LDT theorem statement.
  \item \textbf{Key Mathlib inputs:} elementary arithmetic on nonnegative real
  error parameters.  The project-local inputs are the zero-\(k\) simplification
  of the final error and the degree-zero evaluation fact for polynomial-valued
  measurements.
\end{itemize}

\section{Key theorem forms}

Let \(\sigma\) be the projective strategy in the low individual degree test.
Write \(m\) for the number of variables, \(q\) for the field size, \(d\) for
the individual degree bound, \(k\) for the sampling parameter, and \(\eps\) for
the test-passing error, so that the strategy is assumed to pass the test with
probability at least \(1-\eps\).  The final consistency error in the source
theorem is denoted by \(\nu\).  With this notation there are three relevant
theorem forms.
\begin{enumerate}
  \item \textbf{Printed source form.}  The theorem
  \path{thm:main-formal} assumes \(k\ge md\) and gives the three final
  consistency conclusions at
  \[
    \nu =
    100000 k^2m^4\left(\eps^{1/40000}+(d/q)^{1/40000}
      + e^{-k/(2560000m^2)}\right).
  \]
  \item \textbf{Current checked nonzero form.}  The present formal route proves
  the two-space source conclusion after the source-range work has supplied the
  role-register construction and after the scalar-cascade boundary \(0<k\) is
  available.
  \item \textbf{Zero-sampling boundary.}  If \(d=0\) and \(k=0\), the printed
  hypothesis \(k\ge md\) is satisfied, but the displayed error is exactly
  \(0\).  The theorem then demands exact consistency with degree-zero
  polynomial measurements.
\end{enumerate}

\section{The source assertion}

The source theorem in \cite{Ji2020LowIndividualDegree} states the final
soundness theorem for an integer \(k\ge md\).%
\footnote{Tracked in \ghissue{422}.  The relevant source passage is
\path{references/ldt-paper/test_definition.tex}, lines~180--202.}
This condition admits the boundary case \(k=0\) exactly when \(d=0\).  At that
boundary the printed error parameter is
\[
  100000\cdot 0^2\cdot m^4
  \left(\eps^{1/40000}+(0/q)^{1/40000}
    + e^{-0/(2560000m^2)}\right)=0.
\]
Thus the conclusion is no longer an approximate consistency statement.  It is
an exact consistency statement.

This is stronger than what the test-passing hypothesis implies.  If
\(\eps\) is allowed to be large, the hypothesis that the strategy passes with
probability at least \(1-\eps\) imposes no exact agreement condition on the
point measurements.  In degree zero, however, every polynomial outcome is a
constant function on \(\F_q^m\), so a polynomial-valued measurement evaluates to
the same answer-valued measurement at every point.  A strategy whose point
measurement gives different deterministic answers at two different points can
therefore satisfy the printed hypotheses for a sufficiently large \(\eps\), but
cannot be exactly consistent with a single degree-zero polynomial measurement
at both points.

For instance, take \(m=1\), \(d=0\), and \(k=0\), and let \(q\ge 2\).  On a
one-dimensional classical Hilbert space, let the point measurement answer
deterministically with value \(0\) at one field point and value \(1\) at another.
With \(\eps=1\), the passing hypothesis is the vacuous inequality that the
acceptance probability is at least \(0\).  But any degree-zero polynomial
measurement has a constant answer on \(\F_q\), and so cannot agree exactly with
both point measurements.  Thus the zero value of \(\nu\) would force a
conclusion that this strategy does not satisfy.

This is not a defect of the final scalar cascade.  The cascade simply records
the printed factor \(k^2\).  At \(k=0\) that factor annihilates all error terms,
including the term involving \(\eps\).  Consequently the advertised conclusion
is exact at \(k=0\) for every value of the test-passing parameter \(\eps\); the
choice \(\eps=1\) in the example above is used only to make the printed
hypothesis immediate.  The obstruction itself is algebraic: a degree-zero
polynomial measurement is constant on the field, whereas the point measurements
in the example are not constant as functions of the queried point.

\section{Comparison with the blueprint and formalization}

The formalization now treats this boundary as a source-statement correction.
The corrected Lean declaration
\leanid{MIPStarRE.LDT.Test.mainFormal\_sourceStatement} assumes \(0<k\), in
addition to the separately confirmed large-\(k\) condition recorded in
\path{docs/paper-gaps/issue-906-main-formal-k-bound.tex}.  Under this corrected
statement the source-boundary reduction and its small-error branch are proved by
the two-space role-register route and the scalar-cascade absorption theorem.

The scalar equality \leanid{mainFormalError\_zero\_k} remains as the formal
calculation explaining why the excluded boundary is singular: the advertised
error is zero at \(k=0\).  Thus the zero-sampling boundary is not another
request for a package or bridge input; it is a statement-level correction.

\section{Conclusion}

The printed hypothesis \(k\ge md\) is too weak at \(d=0,k=0\) for the displayed
final error.  At that point the error parameter is exactly zero, while the
test-passing hypothesis with arbitrary \(\eps\) does not force exact agreement
with a constant polynomial measurement.  A source-faithful formal theorem must
therefore either include a justified nonzero sampling condition, such as
\(0<k\), or replace the zero-boundary conclusion by a statement whose error
does not vanish at \(k=0\).  The present formalization takes the first route:
the corrected source theorem assumes \(0<k\), and the former zero-boundary
obligation has been removed rather than discharged by further packaging.

\bibliographystyle{alpha}
\bibliography{references}

\end{document}
